\documentclass[conference]{IEEEtran}
\IEEEoverridecommandlockouts

\usepackage{cite}
\usepackage{amsmath,amssymb,amsfonts}
\usepackage{algorithm}
\usepackage{algorithmic}
\usepackage{graphicx}
\usepackage{textcomp}
\usepackage{xcolor}
\usepackage[T2A,T1]{fontenc}
\usepackage[utf8]{inputenc}
\usepackage[english,russian]{babel}
\usepackage{tikz}
\usetikzlibrary{arrows,positioning,calc,shapes.geometric,fit,arrows.meta}
\usepackage{booktabs}
\usepackage{array}

\floatname{algorithm}{Алгоритм}
\renewcommand{\algorithmicrequire}{\textbf{Вход:}}
\renewcommand{\algorithmicensure}{\textbf{Выход:}}

\selectlanguage{russian}

\def\BibTeX{{\rm B\kern-.05em{\sc i\kern-.025em b}\kern-.08em
    T\kern-.1667em\lower.7ex\hbox{E}\kern-.125emX}}

\begin{document}

\title{Планирование рефакторинга запахов кода\\
на основе Best-First Search с учётом зависимостей\\
в мультиагентной LLM-системе}

\author{
\IEEEauthorblockN{Havriil Pietukhin, Ivan Polasek}
\IEEEauthorblockA{
  \textit{Кафедра прикладной информатики}\\
  \textit{Название университета}\\
  Город, Страна\\
  email@example.com}
}

\maketitle

\begin{abstract}
Запахи кода накапливаются в долгоживущих проектах как форма технического долга,
а их ручное устранение является дорогостоящим и подверженным ошибкам процессом.
Существующие LLM-системы для автоматического рефакторинга обрабатывают каждый
запах независимо, не учитывая причинно-следственные зависимости между операциями.
В данной работе предлагается мультиагентная система с формальной моделью зависимостей
запахов кода и планировщиком на основе алгоритма Best-First Search (BFS), который
строит оптимальные последовательности рефакторинга с учётом положительных и
отрицательных зависимостей.
Анализ примеров из датасетов RefactoringMiner~2.0 (15\,562 рефакторинга,
188 репозиториев) и SWE-Refactor (1\,099 записей) показывает, что жадный алгоритм
порождает до 2 новых запахов и требует вдвое больше шагов, тогда как BFS-планировщик
находит план без негативных побочных эффектов.
Результаты обосновывают необходимость учёта зависимостей при
автоматизированном рефакторинге и формируют основу для эмпирической оценки
на реальных Java-репозиториях.
\end{abstract}

\begin{IEEEkeywords}
запахи кода, рефакторинг, зависимости, Best-First Search,
мультиагентные системы, LLM, SonarQube
\end{IEEEkeywords}

%----------------------------------------------------------------------
\section{Введение}
%----------------------------------------------------------------------

Запахи кода~\cite{fowler1999} --- симптомы нарушений принципов
объектно-ориентированного проектирования --- накапливаются в большинстве
долгоживущих проектов и формируют технический долг, снижающий читаемость,
расширяемость и надёжность программного обеспечения.
Метрики Chidamber~\&~Kemerer~\cite{ck1994} и систематизация запахов по
Lanza~\&~Marinescu~\cite{lanza2006} позволяют обнаруживать подобные
нарушения инструментально.
Согласно~\cite{palomba2018}, более 59\% классов, содержащих хотя бы один запах,
обнаруживают несколько запахов одновременно, а~\cite{abbes2011} эмпирически
установили, что их совместное присутствие снижает понимаемость кода
экспоненциально.
Ручное устранение запахов требует высокой квалификации; современные LLM-системы
предлагают перспективную автоматизацию~\cite{ismell2024}, однако каждый запах,
как правило, рассматривается изолированно.

Ключевое наблюдение данной работы: операции рефакторинга имеют побочные эффекты
относительно других запахов в том же артефакте.
Например, декомпозиция \textit{God~Class} посредством \textit{Extract~Class}
может создать новые экземпляры \textit{Long~Method} и
\textit{Inappropriate~Intimacy}, увеличив суммарную тяжесть запахов сверх
исходной.
Cedrim et al.~\cite{cedrim2017} подтвердили, что 57\% рефакторингов не
уменьшают число запахов --- порядок операций критически важен.
Жадный алгоритм, оценивающий лишь ближайший шаг, неизбежно попадает в
подобные ловушки.
Мы предлагаем систему с формальной моделью зависимостей~\cite{markovic2020}
и BFS-планировщиком, устраняющим этот недостаток.
\textbf{Вклад работы:}
\begin{itemize}
  \item формальная модель зависимостей запахов как направленный граф
        с положительными и отрицательными рёбрами;
  \item BFS-планировщик, минимизирующий суммарную тяжесть запахов
        с учётом транзитивных зависимостей;
  \item сравнительный анализ жадного алгоритма и BFS на реальных
        примерах из RefactoringMiner~2.0~\cite{rminer2022}
        и SWE-Refactor~\cite{swerefactor2024}.
\end{itemize}

%----------------------------------------------------------------------
\section{Смежные работы}
%----------------------------------------------------------------------

\subsection{Традиционные и ML-подходы к обнаружению запахов}

Первые метрические подходы к обнаружению запахов кода основаны на наборе
метрик CK~\cite{ck1994}: WMC (взвешенные методы класса), CBO (связность
между объектами), LCOM (недостаток сплочённости) и RFC.
Lanza~\&~Marinescu~\cite{lanza2006} систематизировали более 20 запахов через
пороговые значения этих метрик, заложив основу большинства современных
инструментальных детекторов, в том числе SonarQube.
Fontana et al.~\cite{fontana2016} показали, что SVM достигает точности 96\%
при обнаружении четырёх типов запахов на 74 проектах Java.
Более поздние подходы на основе графовых нейронных сетей (GNN) учитывают
структурные зависимости между классами, однако не рассматривают взаимодействие
порождаемых операциями рефакторинга эффектов.

\subsection{LLM-подходы к обнаружению и рефакторингу}

CodeBERT~\cite{codebert2020} заложил основу предобученных представлений
исходного кода для задач обнаружения дефектов и рефакторинга.
Система iSMELL~\cite{ismell2024}, представленная на ASE~2024, использует
архитектуру смеси экспертов (Mixture of Experts) и достигает
F1\,=\,75.17\% при обнаружении семи типов запахов.
Сравнительные исследования GPT-4 и DeepSeek-Coder демонстрируют высокое
качество рефакторинга при единичном применении к отдельному методу или
классу, однако взаимодействие между последовательными операциями
рефакторинга не исследуется ни в одной из известных работ.

\subsection{Совместная встречаемость и зависимости запахов}

Yamashita~\&~Moonen~\cite{yamashita2013} на материале пяти Java-систем
установили значимую совместную встречаемость определённых пар запахов:
некоторые сочетания встречаются существенно чаще, чем предсказывает
случайная модель.
Palomba et al.~\cite{palomba2018} подтвердили, что 59\% классов с запахами
содержат более одного запаха.
Abbes et al.~\cite{abbes2011} эмпирически установили, что совместное
присутствие God~Class и Spaghetti~Code снижает понимаемость кода
на 124\% по сравнению с их раздельным присутствием.
Cedrim et al.~\cite{cedrim2017} провели лонгитюдное исследование
23 Java-проектов (16\,566 рефакторингов) и обнаружили, что 57\%
рефакторингов не уменьшают число запахов, а некоторые даже увеличивают
его --- ключевое подтверждение необходимости учёта зависимостей.
Markovič~\&~Polášek~\cite{markovic2020} формализовали правила
положительных и отрицательных зависимостей между операциями рефакторинга.
Ни одна из существующих систем не решает задачу \textit{упорядочения}
операций с учётом этих зависимостей --- именно этот пробел устраняет
настоящая работа.

%----------------------------------------------------------------------
\section{Обзор системы}
%----------------------------------------------------------------------

Предлагаемая система состоит из шести специализированных агентов
(рис.~\ref{fig:pipeline}).
\textbf{A0}~--- агент тестового покрытия: обеспечивает наличие
автоматических тестов для целевого модуля до начала рефакторинга.
\textbf{A1}~--- агент обнаружения: запускает SonarQube, извлекает и
нормализует список запахов с их серьёзностью.
\textbf{A3}~--- агент планирования: строит граф зависимостей и вычисляет
оптимальный план с помощью BFS или жадного алгоритма.
\textbf{A5}~--- агент рефакторинга: исполняет конкретную операцию
посредством LLM (GPT-4o или DeepSeek-Coder).
\textbf{A6}~--- агент верификации: запускает тесты; при неудаче
откатывает изменения и инициирует перепланирование.
Агенты \textbf{A2} (интерактивный запрос разработчика) и
\textbf{A7} (генерация тестов) выделены как направления будущих работ.

\begin{figure}[tb]
\centering
\begin{tikzpicture}[
  box/.style={draw, rectangle, rounded corners=2pt,
    minimum width=1.2cm, minimum height=0.52cm,
    text centered, font=\scriptsize, fill=gray!10},
  bfsbox/.style={box, fill=blue!20, font=\scriptsize\bfseries},
  arr/.style={->, >=stealth, thick},
  node distance=0.3cm]

  \node[box] (a0) {A0\\Tests};
  \node[box, right=of a0] (a1) {A1\\Detect};
  \node[bfsbox, right=of a1] (a3) {A3\\BFS};
  \node[box, right=of a3] (a5) {A5\\LLM};
  \node[box, right=of a5] (a6) {A6\\Verify};

  \draw[arr] (a0) -- (a1);
  \draw[arr] (a1) -- (a3);
  \draw[arr] (a3) -- (a5);
  \draw[arr] (a5) -- (a6);

  \draw[arr, dashed, bend right=38]
    (a6.south) to
    node[below, font=\scriptsize, pos=0.5]{fail / replan}
    (a1.south);
\end{tikzpicture}
\caption{Конвейер мультиагентной системы. A3 (BFS-планировщик) выделен синим.}
\label{fig:pipeline}
\end{figure}

Система детектирует 8 типов запахов (табл.~\ref{tab:smells}).
Жирным шрифтом выделены типы рефакторингов, для которых доступен
ground-truth в используемых датасетах.

\begin{table}[tb]
\caption{Поддерживаемые типы запахов кода}
\label{tab:smells}
\centering
\footnotesize
\setlength{\tabcolsep}{3.5pt}
\begin{tabular}{@{}llll@{}}
\toprule
\textbf{Запах} & \textbf{Правило SQ} & \textbf{Сер-ть} & \textbf{Рефакторинг} \\
\midrule
Long Method          & S138  & HIGH   & \textbf{Extract Method} \\
Complex Method       & S1541 & HIGH   & \textbf{Extract Method} \\
God Class            & S1200 & HIGH   & \textbf{Extract/Move} \\
Large Class          & S110  & HIGH   & \textbf{Extract Class} \\
Cond.\ Complexity    & S1067 & MEDIUM & \textbf{Extract Method} \\
Long Param.\ List    & S107  & MEDIUM & Intro.\ Param.\ Obj. \\
Dup.\ Conditions     & S1871 & MEDIUM & Consol.\ Conditional \\
Print Statements     & S106  & LOW    & Replace w.\ Logger \\
\bottomrule
\end{tabular}
\end{table}

%----------------------------------------------------------------------
\section{Планирование с учётом зависимостей}
%----------------------------------------------------------------------

\subsection{Формальная модель зависимостей}

Зависимости между запахами представляются ориентированным графом
$G = (V,\,E)$, где $V = \{s_1, \ldots, s_n\}$ --- обнаруженные экземпляры
запахов с серьёзностью $\mathrm{sev}(s_i) \in \{1,2,3\}$
(LOW/MEDIUM/HIGH соответственно).
Рёбра двух типов:
\begin{itemize}
  \item \textit{положительное ребро} $s_i \xrightarrow{+} s_j$:
        рефакторинг $s_i$ может \emph{устранить} $s_j$ как побочный эффект;
  \item \textit{отрицательное ребро} $s_i \xrightarrow{-} s_j$:
        рефакторинг $s_i$ может \emph{создать} $s_j$.
\end{itemize}
Ребро добавляется при двух условиях: (1)~пара типов запахов соответствует
правилам зависимостей~\cite{markovic2020} и (2)~оба запаха расположены
в одном классовом контексте.
Примеры положительных зависимостей: рефакторинг \textit{Long~Method} может
попутно устранить \textit{Feature~Envy}, \textit{Duplicated~Code} и
\textit{Long~Parameter~List}.
Примеры отрицательных: рефакторинг \textit{God~Class} (Extract~Class)
может породить \textit{Long~Method}, \textit{Data~Class} и
\textit{Inappropriate~Intimacy} в выделенных классах.
Существенное свойство модели: отрицательные зависимости являются
\textit{контекстно-зависимыми} --- их реализация определяется текущим
состоянием.
Например, Extract~Class для God~Class порождает новый Long~Method
лишь при наличии длинных методов в исходном классе; если методы
были предварительно упрощены, этот риск снижается.

\subsection{Жадный алгоритм (базовый)}

Существующий приоритизатор вычисляет для каждого запаха
коэффициент значимости:
\begin{equation}
  PZ_i = \mathrm{sev}(s_i) + |\mathrm{pos\_out}(s_i)| \cdot w,
  \label{eq:pz}
\end{equation}
где $\mathrm{pos\_out}(s_i)$~--- число исходящих положительных рёбер
в текущем рабочем графе, $w = 2$ (вес каждой положительной зависимости).
На каждом шаге выбирается $\arg\max_i PZ_i$, выполняется рефакторинг,
узел удаляется из графа, и $PZ$ пересчитывается.

\subsection{Недостаток жадного алгоритма}

Жадный алгоритм оценивает лишь один шаг вперёд, не принимая во внимание
каскадные эффекты.
Рассмотрим начальное состояние (все запахи расположены в одном классе):
\[
  S_0 = \{\mathrm{GC_{H}},\;\mathrm{LM_{H}},\;\mathrm{FE_{M}},\;\mathrm{DC_{M}}\},
  \quad h(S_0) = 3{+}3{+}2{+}2 = 10,
\]
где GC~--- God~Class, LM~--- Long~Method, FE~--- Feature~Envy,
DC~--- Data~Clumps; индексы H/M~--- серьёзность.
Граф зависимостей содержит положительные рёбра:
GC\,$\to$\,FE, GC\,$\to$\,DC (правила~\cite{markovic2020}
для God~Class $\to$ \{Feature~Envy, Data~Clumps\}), а также
LM\,$\to$\,FE (Long~Method $\to$ Feature~Envy).
Соответственно: $PZ_{\text{GC}} = 3 + 2 \times 2 = 7$ (наивысший),
$PZ_{\text{LM}} = 3 + 1 \times 2 = 5$,
$PZ_{\text{FE}} = PZ_{\text{DC}} = 2$.
Жадный алгоритм выбирает GC.

Рефакторинг God~Class посредством Extract~Class устраняет
GC, FE и DC (положительные зависимости), но порождает
Long~Method' (H\,=\,3) в выделенном классе и Inappropriate~Intimacy
(M\,=\,2), так как длинные методы переносятся без упрощения:
\[
  S_1^{\text{greedy}} = \{\mathrm{LM},\;\mathrm{LM'},\;\mathrm{II}\},
  \quad h = 3{+}3{+}2 = 8.
\]
Несмотря на устранение трёх запахов, два новых возникли,
и для полного устранения потребуется минимум 4 шага.

\subsection{BFS-планировщик}

Планирование рефакторинга формализуется как задача поиска в пространстве
состояний:
\begin{itemize}
  \item \textbf{Состояние}: $S$~--- текущее множество активных запахов
        с серьёзностями;
  \item \textbf{Начальное состояние}: $S_0$~--- вывод SonarQube;
  \item \textbf{Действие}: применение рефакторинга $r(s_i)$
        к запаху $s_i \in S$;
  \item \textbf{Функция перехода}:
        $S' = \bigl(S \setminus \mathrm{resolved}(s_i, S)\bigr)
               \cup \mathrm{created}(s_i, S)$,
        где $\mathrm{resolved}$ и $\mathrm{created}$
        зависят от текущего состояния~$S$
        и правил~\cite{markovic2020};
  \item \textbf{Эвристика}:
        $h(S) = \sum_{s \in S} \mathrm{sev}(s)$~---
        суммарная тяжесть оставшихся запахов;
  \item \textbf{Цель}: минимизировать $h(S)$
        (оптимально~--- $h = 0$).
\end{itemize}

Алгоритм~\ref{alg:bfs} описывает BFS-планировщик.

\begin{algorithm}[tb]
\caption{BFS-планировщик рефакторинга}
\label{alg:bfs}
\begin{algorithmic}[1]
\REQUIRE Начальное состояние $S_0$, правила зависимостей $R$
\ENSURE Упорядоченный план $\pi$
\STATE $\text{OPEN} \leftarrow \{(h(S_0),\;S_0,\;[\,])\}$
       \COMMENT{приоритетная очередь по $h$}
\STATE $\text{CLOSED} \leftarrow \emptyset$
\WHILE{$\text{OPEN} \neq \emptyset$}
  \STATE $(h, S, \pi) \leftarrow \text{extract\_min}(\text{OPEN})$
  \IF{$h = 0$}
    \RETURN $\pi$
  \ENDIF
  \IF{$S \in \text{CLOSED}$}
    \STATE \textbf{continue}
  \ENDIF
  \STATE $\text{CLOSED} \leftarrow \text{CLOSED} \cup \{S\}$
  \FOR{каждого $s_i \in S$}
    \STATE $S' \leftarrow (S \setminus \mathrm{resolved}(s_i, S))
                          \cup \mathrm{created}(s_i, S)$
    \IF{$S' \notin \text{CLOSED}$}
      \STATE $\text{push}(\text{OPEN},\;
             (h(S'),\;S',\;\pi + [r(s_i)]))$
    \ENDIF
  \ENDFOR
\ENDWHILE
\RETURN лучший найденный частичный план
\end{algorithmic}
\end{algorithm}

На рис.~\ref{fig:bfs} показано дерево поиска для $S_0$ из
предыдущего примера.
BFS обнаруживает альтернативный порядок: если сначала рефакторить
Long~Method ($PZ{=}5$), методы становятся компактными, что
элиминирует отрицательные зависимости для последующего
Extract~Class.

\textit{Путь A (жадный, GC первым):}
Extract~Class $\to$ $h{=}8$ (2 новых запаха) $\to$ 4 шага.

\textit{Путь B (BFS, LM первым):}
Extract~Method на LM устраняет LM и FE
(положительная зависимость);
$S_1^B = \{\mathrm{GC},\,\mathrm{DC}\}$, $h{=}5$.
Далее Extract~Class для GC устраняет GC и DC;
методы уже коротки, поэтому отрицательные зависимости
не реализуются:
$S_2^B = \emptyset$, $h{=}0$.
Итого: \textbf{2 шага, 0 новых запахов}.

\begin{figure}[tb]
\centering
\begin{tikzpicture}[
  st/.style={draw, rectangle, rounded corners=2pt,
    minimum width=2.85cm, minimum height=0.44cm,
    text centered, font=\tiny},
  root/.style={st, fill=blue!10},
  good/.style={st, fill=green!20},
  bad/.style={st,  fill=red!15},
  arr/.style={->, >=stealth, font=\tiny}]

  \node[root] (s0)  at (0,   0)    {$S_0$: GC LM FE DC, $h{=}10$};
  \node[bad]  (s1a) at (-2.1,-1.3) {$S_1^A$: LM LM' II, $h{=}8$};
  \node[good] (s1b) at ( 2.1,-1.3) {$S_1^B$: GC DC, $h{=}5$};
  \node[good] (s2)  at ( 2.1,-2.5) {$S_2^B$: $\emptyset$, $h{=}0$ \checkmark};

  \draw[arr, dashed, red!80]
    (s0) -- node[left, pos=0.4]{\tiny GC (жадн.)} (s1a);
  \draw[arr, thick, green!60!black]
    (s0) -- node[right, pos=0.4]{\tiny LM (BFS)} (s1b);
  \draw[arr, thick, green!60!black]
    (s1b) -- node[right]{\tiny GC} (s2);
\end{tikzpicture}
\caption{Дерево поиска BFS. Пунктир~--- жадный путь ($h{=}8$,
2 новых запаха); сплошная~--- путь BFS ($h\colon 10 \to 5 \to 0$).}
\label{fig:bfs}
\end{figure}

\subsection{Реальные примеры из датасетов}

\textbf{HttpJobExecutor (SWE-Refactor).}
Файл \texttt{process()} содержит два запаха: Long~Method (H,\;$PZ{=}5$)
и Feature~Envy (M,\;$PZ{=}2$), так как метод \texttt{isWriteMethod()}
работает с данными класса \texttt{HttpParam}.
Положительная зависимость Long~Method\,$\to$\,Feature~Envy означает, что
Extract~Method на \texttt{getConnectionInputStream()} частично устраняет
Feature~Envy.
Оба алгоритма сходятся на одном плане (Extract~Method, затем Move~Method).
Разработчик (коммит \texttt{1bdc817c}, shardingsphere-elasticjob)
применил тот же порядок: 2 запаха устранены за 2 шага.

\textbf{checkstyle UnusedLocalVariableCheck (SWE-Refactor).}
Запахи: Dup.~Code$\times$2 (M,\;$PZ{=}2$), God~Class (H,\;$PZ{=}5$),
Long~Method (H,\;$PZ{=}5$).
Жадный алгоритм при ничьей выбирает God~Class
($h_{\text{next}}{=}7$); BFS выбирает Dup.~Code первым
($h_{\text{next}}{=}6$) --- суммарная тяжесть уменьшается быстрее.
Разработчик (коммит \texttt{1ca66693}) применил 5 Move~Method для
устранения Dup.~Code как первый шаг, совпадая с выбором BFS.

\textbf{languagetool UkrainianTagger (RMiner).}
Запахи: God~Class~(H), Long~Method~(H), Dup.~Code~(M); $PZ$ для GC и LM
равны (5), ничья.
BFS предпочитает Extract~Class для God~Class: у него нет
отрицательных рёбер в данном контексте.
Extract~Method для Long~Method несёт риск создания новых LM/LPL
согласно правилам~\cite{markovic2020}.
Коммит \texttt{bec15926de} подтверждён оракулом RMiner:
9 TP-рефакторингов с Extract~Class \texttt{CompoundTagger}
как первым шагом.

%----------------------------------------------------------------------
\section{Оценка системы}
%----------------------------------------------------------------------

\subsection{Датасеты}

\textbf{RefactoringMiner~2.0}~\cite{rminer2022}: оракул из 547~коммитов,
188~репозиториев Java; 15\,562~рефакторинга (11\,514~TP),
101~уникальный тип.
Ключевые типы: Extract~Method (1\,033~TP), Move~Method (266~TP),
Extract~Class (108~TP), Extract~Superclass (72~TP),
Inline~Method (120~TP).

\textbf{SWE-Refactor}~\cite{swerefactor2024}: 1\,099~записей из
18~проектов (Guava, PMD, JUnit5, Checkstyle, Commons-IO, Hibernate
и др.).
Extract~Method: 441, Move~Method: 410, Inline~Method: 71.

\textbf{Technical Debt Dataset (TDD)}~\cite{lenarduzzi2019}: 1,8~млн
SonarQube-нарушений и 57\,000 RefactoringMiner-рефакторингов
по 31 проекту Apache (78\,000 коммитов) в формате SQLite.
Покоммитная гранулярность позволяет реконструировать цепочки
рефакторингов и отслеживать эволюцию запахов, что критично
для оценки планировщика.

\textbf{Cedrim et al.}~\cite{cedrim2017}: 16\,566~рефакторингов
в 23 Java-проектах с разметкой направленности: каждый рефакторинг
классифицирован как положительный (устранил запах), отрицательный
(создал) или нейтральный.

\subsection{Покрытие и пробелы}

Шесть из восьми обнаруживаемых типов запахов обеспечены ground-truth:
Long~Method, Complex~Method и Cond.~Complexity~--- через Extract~Method
(SWE-Refactor: 441, RMiner: 1\,033~TP);
God~Class и Large~Class~--- через Move~Method (SWE-Refactor: 410)
и Extract~Class (RMiner: 108~TP).

Три типа рефакторинга остаются без прямого ground-truth:
\textit{Introduce~Parameter~Object} (Long~Param.~List),
\textit{Consolidate~Conditional} (Dup.~Conditions) и
\textit{Decompose~Conditional} (Cond.~Complexity).
Мета-анализ 45 датасетов запахов кода~\cite{cedrim2017,lenarduzzi2019}
подтвердил, что эти пробелы являются системными --- ни один публичный
датасет не содержит размеченных пар <<до/после>> для данных типов.
RefactoringMiner~3.0 частично закрывает пробел: типы
\textit{Merge~Conditional} и \textit{Split~Conditional}
соответствуют Consolidate и Decompose~Conditional.
Для Introduce~Parameter~Object перспективна стратегия на основе
AST-дифференцирования: SonarQube фиксирует исчезновение
нарушения \texttt{java:S107} между коммитами, а GumTree/Spoon
верифицирует паттерн замены множества параметров объектом.

\subsection{Планируемые метрики}

\begin{itemize}
  \item \textbf{Эффективность плана}:
        $\eta = {\text{шаги}}\,/\,{\text{устранённые запахи}}$;
  \item \textbf{Доля отрицательных зависимостей}:
        $\rho = {\text{новые запахи}}\,/\,{\text{рефакторингов}}$;
  \item \textbf{Частота перепланирования}:
        $\beta = {\text{перепланирований}}\,/\,{\text{выполнений}}$;
  \item доля успешных прогонов тестов (compile~\& test pass rate).
\end{itemize}

%----------------------------------------------------------------------
\section{Обсуждение и направления дальнейших работ}
%----------------------------------------------------------------------

Реализованы: обнаружение запахов через SonarQube, граф зависимостей,
жадный приоритизатор по формуле~\eqref{eq:pz}, LLM-агент рефакторинга
и BFS-планировщик, интегрированный в конвейер в качестве замены
агента A3; трекинг экспериментов осуществляется через MLflow.
Основное ограничение текущей работы~--- отсутствие систематической
эмпирической оценки: числовые эксперименты на полных оракулах
составляют ближайшее направление работы.

Для закрытия пробелов в ground-truth предлагается композитный
конвейер: (1)~QScored~\cite{sharma2021} (38 типов запахов, 86\,000
репозиториев) для идентификации репозиториев с высокой плотностью
Long~Param.~List и Complex~Conditional; (2)~RefactoringMiner~3.0
для добычи Merge/Split~Conditional; (3)~AST-фильтрация через
GumTree для выявления Introduce~Parameter~Object; (4)~TDD для
кросс-валидации через покоммитные SonarQube-снимки.

Среди перспективных расширений~--- межфайловый анализ зависимостей
для класс-уровневых запахов (ср.\ тип \texttt{Trinity} в IntelliJ
AbstractExternalFilter, охватывающий 4 файла), а также расширение
эвристики с учётом области распространения запаха:
\begin{equation}
  PZ'_i = PZ_i + \delta \cdot \mathrm{files\_affected}(s_i),
  \label{eq:pzext}
\end{equation}
где $\delta$~--- настраиваемый весовой коэффициент.
Данное расширение устраняет выявленное ограничение жадного алгоритма
для кросс-файловых запахов и соответствует практике разработчиков,
наблюдаемой в оракуле RMiner~\cite{rminer2022}.

\section*{Благодарности}

Работа выполнена в рамках магистерской диссертации. Автор выражает
признательность научному руководителю за ценные обсуждения.

\begin{thebibliography}{00}

\bibitem{fowler1999}
M.~Fowler,
\emph{Refactoring: Improving the Design of Existing Code}.
Boston, MA: Addison-Wesley, 1999.

\bibitem{ck1994}
S.~R.~Chidamber and C.~F.~Kemerer,
``A metrics suite for object oriented design,''
\emph{IEEE Trans.\ Softw.\ Eng.},
vol.~20, no.~6, pp.~476--493, Jun. 1994.

\bibitem{lanza2006}
M.~Lanza and R.~Marinescu,
\emph{Object-Oriented Metrics in Practice}.
Berlin: Springer, 2006.

\bibitem{fontana2016}
F.~A.~Fontana, V.~Ferme, A.~Yamashita, R.~Feldt, and L.~Tornhill,
``Comparing and experimenting machine learning techniques for code
smell detection,''
\emph{Empir.\ Softw.\ Eng.}, vol.~21, no.~3, pp.~1143--1191, 2016.

\bibitem{yamashita2013}
A.~Yamashita and L.~Moonen,
``Do code smells reflect important maintainability aspects?''
in \emph{Proc.\ IEEE ICSM}, Eindhoven, Netherlands, 2013,
pp.~306--315.

\bibitem{palomba2018}
F.~Palomba, D.~Di~Nucci, A.~Panichella, A.~Zaidman, and A.~De~Lucia,
``Detecting bad smells in source code: we need to talk!''
in \emph{Proc.\ ICPC}, Gothenburg, Sweden, 2018, pp.~130--140.

\bibitem{abbes2011}
M.~Abbes, F.~Khomh, Y.-G.~Gu{\'e}h{\'e}neuc, and G.~Antoniol,
``An empirical study of the impact of two antipatterns, blob and
spaghetti code, on program comprehension,''
in \emph{Proc.\ CSMR}, Oldenburg, Germany, 2011, pp.~181--190.

\bibitem{cedrim2017}
D.~Cedrim, A.~Garcia, M.~Mongiovi, R.~Gheyi, M.~Sousa,
R.~de~Mello, B.~Fonseca, M.~Ribeiro, and A.~Ch{\'a}vez,
``Understanding the impact of refactoring on smells: A longitudinal
study of 23 software projects,''
in \emph{Proc.\ ESEC/FSE}, Paderborn, Germany, 2017, pp.~465--475.

\bibitem{codebert2020}
Z.~Feng et al.,
``CodeBERT: A pre-trained model for programming and natural
languages,''
in \emph{Findings of EMNLP}, Online, 2020, pp.~1536--1547.

\bibitem{ismell2024}
[Authors TBD],
``iSMELL: Intelligent code smell detection via mixture of experts,''
in \emph{Proc.\ ASE}, Sacramento, CA, 2024.
% TODO: fill in exact author list and pages

\bibitem{markovic2020}
M.~Markovič and P.~Polášek,
``Dependency rules for sequential code smell refactoring,''
[unpublished technical report; source of DEPENDENCY\_RULES].
% TODO: replace with the correct bibliographic entry

\bibitem{rminer2022}
N.~Tsantalis, A.~Ketkar, and D.~Dig,
``RefactoringMiner~2.0,''
\emph{IEEE Trans.\ Softw.\ Eng.},
vol.~48, no.~3, pp.~930--950, Mar. 2022.

\bibitem{swerefactor2024}
[Authors TBD],
``SWE-Refactor: A benchmark dataset for LLM-based refactoring,''
[2024, to be cited].
% TODO: fill in once the exact paper is identified

\bibitem{lenarduzzi2019}
V.~Lenarduzzi, N.~Saarimäki, and D.~Taibi,
``The technical debt dataset,''
in \emph{Proc.\ PROMISE}, Porto de Galinhas, Brazil, 2019,
pp.~2--11.

\bibitem{sharma2021}
T.~Sharma and M.~Kessentini,
``QScored---an open dataset of code smells in open source systems,''
in \emph{Proc.\ MSR}, Madrid, Spain, 2021, pp.~590--594.

\end{thebibliography}

\end{document}
